\documentclass[10pt,b5paper,openany]{jsbook}

\usepackage[dvipdfmx]{graphicx}
\usepackage[dvipdfmx]{color}
\usepackage{listings}
\usepackage{inconsolata}
\usepackage{url}
\usepackage{titlesec}
\usepackage{setspace}
\usepackage{wrapfig}

\lstset{basicstyle={\footnotesize\ttfamily}}
\lstset{breaklines=true}
\lstset{lineskip=-3pt}
\lstset{aboveskip=16pt}
\lstset{belowskip=10pt}

\renewcommand{\kanjifamilydefault}{\gtdefault}
\renewcommand{\familydefault}{\sfdefault}
\renewcommand{\contentsname}{}
\renewcommand{\baselinestretch}{0.9}

\definecolor{gray75}{gray}{0.75}
\titleformat{\chapter}[hang]{\Huge\bfseries}{\thechapter\textcolor{gray75}{｜}}{0pt}{\Huge\bfseries}

\begin{document}

{\Huge OpenStackの本 --- 実践編}

\vspace*{10mm}
\begin{minipage}{0.5\textwidth}
	{\scriptsize 毎度おなじみOpenStackの本ですが、今回でついに4冊目になりました。絞ったぞうきんをさらに絞るような感が出てきましたが、今回は実践編ということにして、ログに関することを書くことにしました。それだけではページが足りないので、ついでにOpenStackのサブプロジェクトまとめのようなものを後半に持ってきました。}
\end{minipage}

\vspace*{-1truein}
\begin{minipage}{\textwidth}
	\tableofcontents
\end{minipage}

\thispagestyle{empty}

\documentclass[9pt,b5paper,tombo,openany]{jsbook}

\usepackage[dvipdfmx]{graphicx}
\usepackage{listings}
\usepackage{inconsolata}
\usepackage{url}
\usepackage{titlesec}
\usepackage[dvipdfmx]{color}
\usepackage{setspace}
\usepackage{wrapfig}

\lstset{basicstyle={\small\ttfamily}}
\lstset{frame=single}
\lstset{breaklines=true}
\lstset{lineskip=-1.5pt}
\lstset{framesep=6pt}

\renewcommand{\kanjifamilydefault}{\gtdefault}
\renewcommand{\familydefault}{\sfdefault}
\renewcommand{\contentsname}{}
\renewcommand{\baselinestretch}{1.1}

\definecolor{gray75}{gray}{0.75}
\titleformat{\chapter}[hang]{\Huge\bfseries}{\thechapter\textcolor{gray75}{｜}}{0pt}{\Huge\bfseries}

\begin{document}

\noindent
{\Huge OpenStackの本 番外編}

\vspace*{-1in}
\begin{minipage}{0.4\paperwidth}
	\tableofcontents
\end{minipage}

\vspace*{1in}
\begin{minipage}{0.4\paperwidth}
	\begin{spacing}{0.75}
	\end{spacing}
\end{minipage}

\thispagestyle{empty}

\chapter{OpenStackアップグレード超入門}

\setcounter{page}{1}

\marginpar{\includegraphics[width=30mm]{mini1.png}}本章は、 Qiita に投稿した OpenStack Advent Calendar 2015 の 12 月 3 日の記事である「 OpenStack をアップグレードしたら心臓止まりかけた話」という物語を改めて書き記した章である(ただし、改訂版というわけではない)。

\section{ログについて考えてみる}
\section{Fluentdについて}
\section{Kafkaについて}
\section{InfluxDBについて}
\section{OpenStack運用のためにチェックしたほうが良いログについて}
\section{KVMからVMのメトリクスを取得する方法}




\chapter{あとがき}

\setstretch{0.96}

\section*{こじろー}

\marginpar{\includegraphics[width=30mm]{mini3.png}}

\section*{まっきーさん}

\marginpar{\includegraphics[width=30mm]{mini3.png}}

\noindent
代わって檻には一匹の精悍な豹が入れられた。

\noindent
（西岡兄姉（2010）『カフカ・断食芸人』ヴィレッジブックス）

\vspace*{\stretch{1}}
\setstretch{1.0}
\begin{minipage}{0.5\paperwidth}
	著者：こじろー・まっきー

	挿絵：かとう

	発行：2016年12月31日

	印刷：POPLS (\url{http://www.inv.co.jp/~popls/})
\end{minipage}

\end{document}


\chapter{Big-Tent解説}
OpenStackはそれぞれのコンポーネントが疎結合であるため、連携が容易であることが利点の１つです。その結果、さまざまな連携プロジェクトが生まれ、非常にリッチなクラウド環境の構築ができるようになりました。品質を維持するために、APIの仕様はきちんと議論された上で定義されます。それぞれのコンポーネントでバージョンを一致させ、結合テストを行います。が、つまりは増えすぎました。さまざまな派生プロジェクト間での調整を行うことが難しくなり、足並みをそろえるだけで数ヶ月かかるようになっていまいました。そこで、OpenStackのコアとなるコンポーネントを定義し、それらの結合テスト以外はサブプロジェクトとして、リリースを一致させる必要はない、ということにしました。これをOpenStack界隈ではBigTent構想と言っています。

とはいえ、プロジェクトを減らそうという話ではないわけなので、相変わらずたくさんのプロジェクトが乱立しています。そんなわけで、サブプロジェクトの一部を紹介してみることにしました。プロジェクトのリストはOpenStackのWikiでまとめられているのですが、すべて英語なので調べる気が起きていないひともいるかと思います。

中には単一ベンダーが１人で頑張っているだけのプロジェクトもあります。また、機能が重複しているんじゃないかというプロジェクトもあります。ちゃんとWikiが整備されていなかったり説明が無かったりと、BigTentは玉石混交です。

\section{CloudKitty}
\begin{description}
	\item[wiki:] \url{https://wiki.openstack.org/wiki/CloudKitty}
	\item[Document:] \url{http://docs.openstack.org/developer/cloudkitty/}
\end{description}
Rating-as-a-serviceと説明されていますが、課金管理をしてくれます。

Ceilometerは確かにメトリクスを計測してくれるのですが、課金に関してはそのスコープ外でした。課金のもととなるデータ(collector)、課金ポリシー(rating pipeline)、課金情報の保管(storage)、書き出し(writers)から構成されています。複数のコレクターや書き出しフォーマットを選ぶことができます。

\section{Congress}
\begin{description}
	\item[wiki:] \url{https://wiki.openstack.org/wiki/Congress}
\end{description}

\section{Designate}
\begin{description}
	\item[wiki:] \url{https://wiki.openstack.org/wiki/Designate}
\end{description}

\includegraphics[width=0.5\textwidth]{img/logo-designate.pdf}

みんな大好きDNSのサービスです。
\begin{itemize}
	\item ドメインとレコードの管理ができるAPIの提供
	\item マルチテナントのサポート
	\item NovaやNeutronと連携してレコードの自動追加
\end{itemize}
いまのところPowerDNSとBIND9をインストール時点でサポート。プラグイン型式で追加も可能のようです。

\section{Dragonflow}
\begin{description}
	\item[wiki:] \url{https://wiki.openstack.org/wiki/Dragonflow}
	\item[Document:] \url{http://docs.openstack.org/developer/dragonflow/distributed_dragonflow.html}
\end{description}
\includegraphics[width=0.5\textwidth]{img/Df_logo.png}

\section{Freezer}
\begin{description}
	\item[wiki:] \url{https://wiki.openstack.org/wiki/Freezer}
\end{description}

\includegraphics[width=0.5\textwidth]{img/freezer_logo.png}

バックアップサービスです。複数のOSをサポートし(Linux・Windows・OSX・*BSD)、ブロックストレージのバックアップやファイルシステムの差分バックアップを可能にすることが目標です。また、指定時間でのバックアップやジョブの同期など、バックアップスケジューリングもスコープに含みます。

\section{Karbor}
\begin{description}
	\item[wiki:] \url{https://wiki.openstack.org/wiki/Karbor}
\end{description}
旧Smaugです。WikiにはOpenStackデプロイのデータとメタデータを保護するもの、あるいはKeystoneのプロジェクト下にあるものの保護、と書いてありますがこれだけではどうにもピンときません。

例えば以下のようなアプリケーションがあったとします。

\includegraphics[width=\textwidth]{img/Smaug-sample-application.png}

このアプリケーションをシステムまるごと保存し、変更されたことを検知したいとします。このシステムを説明すると
\begin{itemize}
	\item DBインスタンスは1つで、DBネットワークに接続されている
	\item DBネットワークはRouter配下にある
	\item Appインスタンスは1つで、Appネットワークに接続されている
	\item AppネットワークはRouter配下にある
	\item Webインスタンスは2つあり
\end{itemize}
などなど、これらの情報を保存しておくことになります。

Karborではこれを依存グラフとして表現します

\includegraphics[width=\textwidth]{img/Smaug-dependency-graph.png}

\section{Kolla}
\begin{description}
	\item[wiki:] \url{https://wiki.openstack.org/wiki/Kolla}
\end{description}
OpenStackのコンポーネントをDocker化しようというプロジェクトです。OpenStackのAPIサーバーはImmutableなものとして扱うことが容易なので、すでにコンテナを使っている人も多いのではないかと思います。MissionにはOpenStackデプロイの簡素化や素早い増強ができるようにするとあります。

特にカスタマイズしなくても動くようなイメージを目指しているようですが、DockerfileをJinja2のテンプレートになっているのでカスタマイズが可能です。ソースコードからのインストール、RDOやUbuntuのリポジトリを使ったインストールの両方をサポートしています。

\section{Kuryr}
\begin{description}
	\item[wiki:] \url{https://wiki.openstack.org/wiki/Kuryr}
\end{description}
コンテナ用に設計されたネットワークプロジェクトです。もとはNeutronのサブプロジェクトでしたが、分離されました。コンテナに関しては長らくすったもんだを繰り返していましたが、ネットワークはKuryrになんとか落ち着いていきそうです。コンテナ自体はMagnumで落ち着きそうですし、OpenStackとコンテナはなんとか落としどころを見つけたのでしょうか。そもそも、仮想マシン(インスタンス)とコンテナではライフサイクルが違う、ということに無理があったそうで。

ところで、どう発音すればよいのかは謎である。

\section{Searchlight}
\begin{description}
	\item[wiki:] \url{https://wiki.openstack.org/wiki/Searchlight}
\end{description}
確か最初はGlanceのイメージを検索できるように、というような趣旨のプロジェクトだった気がするのですが、LibertyでのデザインサミットでGlanceに限らずOpenStackのリソースすべてを検索対象とする方針の転換がありました。今はサーチサービスという大義のようです。ミッションは「マルチテナンシーのある、スケーラブルなインデクサーとサーチエンジンの提供」です。

OpenStackにかかわるあらゆるもの、例えばインスタンスなどの情報を集めて、Elasticsearchにインデックスし、それを検索可能にするAPIを提供します。Horizonのインスタンス検索機能は貧弱ですし、APIも文字列部分一致程度のリスティング機能しかありませんので、Searchlightは便利だと言えるでしょう。

\section{Senlin}
\begin{description}
	\item[インストールガイド:] \url{http://docs.openstack.org/developer/senlin/install.html}
	\item[開発者ガイド:] \url{http://docs.openstack.org/developer/senlin/developer/index.html}
	\item[wiki:] \url{https://wiki.openstack.org/wiki/Senlin}
\end{description}
インスタンスをクラスタ化することができるAPIを提供します。Novaが作成したインスタンスだけが対象ではなく、Heatのスタックなどもクラスタ化することができます。クラスタ化するにおいて、ポリシーを設定することができ、以下のプラグインがあるようです

\begin{itemize}
	\item 配置 (Placement)
	\item 削除 (Deletion)
	\item スケール (Scaling)
	\item ヘルスチェック (Health)
	\item ロードバランシング
	\item その他
\end{itemize}

\section{Solum}
\begin{description}
	\item[wiki:] \url{https://wiki.openstack.org/wiki/Solum}
\end{description}
Software Development Lifecycle Automation service: 直訳するとソフトウエア開発サイクルの自動化サービス、です。ソースコードのOSイメージへの配置やデプロイ作業を自動化し、OpenStackをアプリケーション開発フローと統合させることがミッションです。

プロジェクトのゴールは開発者目線からとアプリケーションの可搬性、アプリケーションを作成するプログラミング言語に依存しない設計の3つから定められています。

\subsection*{開発者の生産性}
\begin{itemize}
	\item 開発・テスト・ステージング・本番、などの環境を意識できるアプリケーションサイクルの管理ができる
	\item 自動化されたデプロイフローによるCI/CD環境の構築
	\item \verb|git push|
	\item IDE連携
\end{itemize}

\subsection*{アプリケーションの可搬性}
\begin{itemize}
	\item Solumnで動くアプリケーションは、どのベンダーのOpenStackでも動く
	\item パブリッククラウドでもプライベートクラウドでも動く
	\item Novaベースでデザインされているため、Dockerに対応
\end{itemize}

\subsection*{言語・拡張性}
\begin{itemize}
	\item ランタイムをプラグイン形式にすることで、Solumn自体はアプリケーションを作成したのプログラミング言語に依存しない。
	\item アドオン型式なので、独自の拡張が可能
	\item つまり、ベンダー独自の拡張も可能ではある
\end{itemize}

NovaのDokcerを使う、という辺りに少し不安はあるものの、ニーズは大きそうなプロジェクトです。

\section{Vitrage}
\begin{description}
	\item[wiki:] \url{https://wiki.openstack.org/wiki/Vitrage}
\end{description}

\includegraphics{img/Vitrage_logo_finaly.png}

OpenStackのアラートやイベントを可視化したり分析することで、トラブルを未然に防いだり、トラブルシューティングに役立て得ようとするプロジェクトです。

アーキテクチャは以下のようになっています

\includegraphics[width=\textwidth]{img/Vitrage-high_level_architecture2.png}

機能としては

\begin{itemize}
	\item コンピュートホストやネットワーク機器などの物理機器の状態と、その物理機に依存している仮想マシンとの紐付け
	\item 直接OpenStackを監視するのではなく、システムの分析をベースにした状態の変化やアラーティング
	\item アラートやイベントによる問題解決
	\item ダッシュボード (Horizon)
\end{itemize}

とあるネットワークスイッチで障害が発生した場合、どの仮想マシンが影響を受けるか、といった分析が可能になります。これらの情報を持っていることで、物理機器のメンテナンス時の影響を把握することも簡単になりそうです。

\section{Watcher}
\begin{description}
	\item[wiki:] \url{https://wiki.openstack.org/wiki/Watcher}
\end{description}
OpenStackベースのクラウド環境でのリソース制御機構を提供しようというプロジェクトです。システムメトリクスの収集や最適化、アクションプランのプロセスが動作します。データセンターの運用コストやマイグレーションによるシステムパフォーマンスの改善、電源効率の改善なども視野に入っているようです。制御アルゴリズムはデフォルトで入っているもの以外にもプラグイン形式で追加することもできます。OpenStack運用者のためのプロジェクトを言えます。なんともピンと来ませんが、いくつがユースケースが紹介されています。

\begin{itemize}
	\item サーバー筐体の温度や空調、電源管理などの物理的なリソースの運用時に、仮想マシンのホスト間移動や適切なスケジューリングの機構を提供することができるようになる。
	\item 最適化が必要になるしきい値の指定や変更が容易。
	\item テストや開発環境ではリソースを詰めて、本番環境では少し余裕を持たせたスケジューリングをしたいなど、リソースセットの定義をすることができる。
	\item リソースの使用状況を詳細に監視することができる。
\end{itemize}

以下のアーキテクチャで動作しています

\begin{itemize}
	\item 基本的にはメッセージキューからのイベント駆動
	\item イベントにはOpenStackのあらゆるコンポーネントに対応可能。 (Novaのマイグレーション、Keystoneの認証など)
	\item 逆にWatcherからイベントを発行させることも可能
\end{itemize}


\chapter{あとがき}

\section*{こじろー}
このような薄い本を手に取っていただきありがとうございます。こんな本のあとがきなんてもう書くことが何もないので、ついに「同人誌 あとがき」でググってみました。そしたら、最近の傾向としては、そもそもあとがきなんか書かない派のサークルさんが多いらいいです。ふむふむ、と思って自分たちの同人誌の過去のあとがきを読むと、もうヒドすぎて、これはもうあとがきではなく、ただの落書きだなって感じでした。こんな落書きだったら、実は書かないほうが良かったんじゃなかろうかって思ったのですが、そういやページ稼ぎのために無理やり尺を伸ばした文章だったと思い出して、残念な気持ちになりました。

ところで、文章を書くという行為は本当に精神力を使うなってヒシヒシと感じております。この同人誌のための文章を書くにしても、少しずつしか書けなくて、こんなに集中力が続かないものかと、自分でもびっくりしている次第です。これはどうやら歳のせいだと勝手に結論付けて納得しておきます。加えて、OpenStackに対する自分の熱量が低くなったからなんだろうなって感じます。はっきり言ってネタが無くて困っており、次回夏コミ(当確すれば)のネタは全く思いつきません。ですが、実は来年から自分はOpenStackに対する熱量を増やしていかなければいけない仕事をすることになりまして、こんな泣き言みたいなことを言っている場合ではないな、自分を律するようにしている次第でござます。

今は一旦、この本が完成することに満足し、一息ついて、また来年は何をしようか、ゆっくり考えてきたいと思う。

\section*{まっきーさん}
書くことない。

\vspace*{\stretch{1}}
\setstretch{1.0}
\begin{minipage}{0.5\paperwidth}
	\begin{description}
		\item[著者:]こじろー・まっきー
		\item[挿絵:]かとう
		\item[発行:]2016年12月31日
		\item[印刷:]POPLS (\url{http://www.inv.co.jp/~popls/})
	\end{description}
\end{minipage}

\end{document}
